\section{Conclusion}
\label{sec:conclusion}

We compared hand-crafted-feature regression (Phase~1) with deep sequence
models and a hybrid stacking model (Phase~2) for RUL prediction on C-MAPSS FD001.
All models use the same engine split, the same 100 fixed validation points and the true uncapped RUL;
the test set was scored once.

\subsection{Key Takeaways}
% Source: results/phase2/final_results.csv, test_bootstrap.csv, ablation_summary.csv
\begin{enumerate}
  \item \textbf{A single-layer LSTM is the best model, but its advantage is not statistically established.}
    It reaches a validation MAE of $20.14\pm1.46$ (19.41 for the saved seed-42 model) against 25.03 for Random Forest,
    and a test MAE of 16.10 against 18.95.  The paired bootstrap 95\% interval of the test difference is $[-0.14, 5.84]$ cycles, which includes 0.
  \item \textbf{Architecture matters less than expected.}  On the validation-window MAE the LSTM, GRU and Transformer are close; on the
    test set the GRU and CNN are worse than both tree models, and the Transformer is level with them.
  \item \textbf{Small data limits capacity.}  With about 80 independent engines, extra recurrent layers hurt in the ablation,
    early stopping is the regularizer that matters, and capping the training target does not help under uncapped evaluation.
  \item \textbf{The hybrid improves a tree model, not the best deep model.}  Stacking lowers the Random Forest MAE from $25.08$ to $21.72$ on validation,
    but it does not beat the LSTM alone, and it is no better than the plain average of the two.  Its lowest NASA score (4312) comes with
    over-prediction of 84\% of the late test points.
\end{enumerate}

\subsection{Limitations}
\begin{itemize}
  \item FD001 has one operating condition and one fault mode; FD002--FD004 were not tested.
  \item There are only 80 training engines, 20 validation engines and 100 test engines.  Confidence intervals are wide, validation scores are slightly optimistic
    because early stopping and model selection used the validation engines, and the ablation standard deviations rest on three seeds.
  \item No model removes the underprediction for engines with a true RUL above 200 cycles (mean error about $-110$ to $-125$ cycles on the 5 validation points that fall there).
    The test set has no such engine, so the gap cannot be measured on it.
  \item The 30-cycle window of a healthy engine contains little information about its remaining life; this is the most likely cause of the gap.
\end{itemize}

\subsection{Future Work}
\begin{itemize}
  \item Evaluate on FD002--FD004 and across subsets.
  \item Use longer histories for healthy engines (for example variable-length or hierarchical windows), which is the most direct remedy for the high-RUL gap.
  \item Add uncertainty estimates (deep ensembles or MC dropout), since a point prediction hides the risk of late predictions.
  \item Try augmentation by jittering or synthetic trajectories to reduce the dependence on 80 engines.
\end{itemize}
